\documentstyle[%


righttag%


,11pt%


,amssymb%


,russian%               remove in Eng.


,izvru%                 Set izven% in Eng.


]{amsart}





% 





\newcommand{\Close}[1]{{\langle {#1}\rangle}}


\newcommand{\Class}[1]{{\Close{\!\Close{#1}\!}}}


\newcommand{\Sub}{\bold{Sub}}


\newcommand{\tts}[1]{}


%\renewcommand{\baselinestretch}{0.96}





%\hoffset=-15mm%                  For Eng.


\setcounter{page}{53}


\god{2000}


\nomer{7~(458)} %                        Remove in Eng.


%\nomer{Vol.\,44, No.\,7}%               Set in Eng.


%\str{\pageref{begin}--\pageref{end}}%  Set in Eng.


\udk{512}





\author{\it А.П.\,Семигродских}


% NB:   ^^ remove in Eng.


\title{О клонах полиномов над бесконечными полями}





\date{}


%\pagestyle{myheadings}%         Set in Eng.


\pagestyle{plain} %              Remove in Eng.


\begin{document}


%\thispagestyle{plain}%          Set in Eng.


\maketitle


\label{begin}





\section{Введение}





Данная работа является продолжением [1] и 


посвящена задаче классификации части


замкнутых классов полиномов над бесконечными полями.





Введем основные понятия. Пусть дан класс функций, 


отображающих декартовы степени множества $K$ в $K$.


Этот класс называется замкнутым,


если он замкнут относительно суперпозиций.


Более точно, для $m$-местной функции $f$ и 


$n$-местной функции $g$ определим


$r$-местную функцию $h=f*g$ ($r=m+n-1$), положив


$$


h(x_1,\dotsc,x_r)=f(g(x_1,\dotsc,x_n),x_{n+1},\dotsc,x_r)


$$


для всех $x_1,\dotsc,x_r\in K$. 


Далее определим операции $\zeta$, $\tau$, 


$\Delta$ и $\nabla$, полагая при $m>1$


\begin{align*}


(\zeta f)(x_1,\dotsc,x_m)&=f(x_2,\dotsc,x_m,x_1),\\


(\tau f)(x_1,x_2,x_3,\dotsc,x_m)&=f(x_2,x_1,x_3,\dotsc,x_m),\\


(\Delta f)(x_1,\dotsc,x_{m-1})&=f(x_1,x_1,x_2,\dotsc,x_{m-1}),\\


(\nabla f)(x_1,\dotsc,x_{m+1})&=f(x_2,\dotsc,x_{m+1})


\end{align*}


для всех $x_1,\dotsc,x_{m+1}\in K$, и $\zeta f=\tau f=\Delta f=f$ при $m=1$. 


Класс функций $C$ называется замкнутым, 


если он замкнут относительно операций $*$, $\zeta$, $\tau$, 


$\Delta$ и $\nabla$. 


Такой подход, предложенный А.И.\,Мальцевым, 


позволяет рассматривать замкнутые классы как


алгебры относительно введенных операций (см.~[2]).


Замкнутый класс называется клоном, если он содержит все проекции,


т.\,е. функции вида $e^n_i(x_1,\dotsc,x_n)=x_i$, $1\le i\le n$.





Пусть теперь $K$ --- произвольное 


ассоциативное коммутативное кольцо с единицей. Мы рассматриваем


алгебраические полиномы над $K$ от произвольного (конечного) числа


переменных и замкнутые классы, образуемые такими полиномами.


Рассмотрим основные для нас замкнутые классы





$F_K$ --- класс всех полиномов над $K$;





$F^0_K$ --- класс всех полиномов из $F_K$ без свободного члена;





$L_K$ --- класс всех линейных полиномов из $F_K$;





$L^0_K$ --- класс всех линейных полиномов из $L_K$ без свободного члена.





Поставим задачу классификации 


клонов из интервала $[{L^0_K}, {F_K}]$ решетки всех клонов.


Будем понимать это как задачу выяснения решеточной


структуры указанного интервала, 


который изображен на диаграмме~1.


Знаки вопросов на ней означают, что указанные интервалы 


могут содержать замкнутые классы, не совпадающие


с четырьмя отмеченными.


Это зависит от свойств кольца $K$.











В [3] было показано,


что в случае кольца вычетов по любому модулю весь решеточный интервал


$[{L^0_K}, {F_K}]$ есть подпрямое произведение интервалов


$[{L^0_K}, {F^0_K}]$ и


$[{L^0_K}, {L_K}]$. Соответствующее доказательство проходит и для


произвольного кольца $K$. Поэтому исследование интервала $[{L^0_K}, {F_K}]$


в каком-то смысле сводится к изучению интервалов $[{L^0_K}, {F^0_K}]$ и


$[{L^0_K}, {L_K}]$. В $[2]$ для случая кольца вычетов показано, что интервалы


$[{L^0_K}, {L_K}]$ и $[{F^0_K}, {F_K}]$ изоморфны решетке идеалов


кольца $K$ (а значит, и между собой). Доказательство и этого факта


без изменений переносится на произвольное кольцо $K$. 


Поэтому, если $K$ --- поле, то


интервал $[{L^0_K}, {F_K}]$ имеет вид,


изображенный на диаграмме~2.





\begin{picture}(350,110)


%% added 50 units to all X


\put(165,85){\oval(40,50)[bl]}


\put(165,35){\oval(40,50)[tl]}


\put(125,85){\oval(40,50)[br]}


\put(125,35){\oval(40,50)[tr]}





\put(145,60){\oval(40,50)}





\put(140,20){$L^0_K$}


\put(140,90){$F_K$}


\put(105,55){$F^0_K$}


\put(170,55){$L_K$}





\put(165,60){\circle*{5}}


\put(145,35){\circle*{5}}


\put(125,60){\circle*{5}}


\put(145,85){\circle*{5}}





\put(132,70){?}


\put(132,43){?}


\put(142,57){?}


\put(155,43){?}


\put(155,70){?}





\put(110,5){Диаграмма 1}





\put(300,85){\oval(40,50)[bl]}


\put(260,35){\oval(40,50)[tr]}





\put(280,60){\oval(40,50)[bl]}


\put(280,60){\oval(40,50)[tr]}





\put(275,20){$L^0_K$}


\put(275,90){$F_K$}


\put(240,55){$F^0_K$}


\put(305,55){$L_K$}





\put(300,60){\circle*{5}}


\put(280,35){\circle*{5}}


\put(260,60){\circle*{5}}


\put(280,85){\circle*{5}}





\put(280,35){\line(4,5){20}}


\put(280,85){\line(-4,-5){20}}





\put(287,70){?}


\put(267,43){?}





\put(245,5){Диаграмма 2}





\end{picture}





Пусть $C$ --- замкнутый класс


из $[{L^0_K}, F_K]$. Тогда $C$ может быть порожден совокупностью 


всех линейных форм (т.\,е. полиномов из $L^0_K$) и какими-то функциями


${f_1}, {f_2}, \dotsc$ из $F_K$. В этом случае будем писать


$C=\Class{{f_1}, {f_2}, \dotsc }$ и говорить, 


что $C$ порожден функциями ${f_1}, {f_2}, \dotsc$ (не упоминая $L^0_K$). 


Если при этом множество $\{ {f_1}, {f_2}, \dotsc \}$ конечно, 


то класс $C$ будем называть конечнопорожденным.


Всюду далее словосочетание


``замкнутый класс'' означает {\em замкнутый класс из $[{L^0_K},F_K]$},


а ``полином'' или ``многочлен'' --- {\em полином над $K$}. 


Каждый из изучаемых классов содержит $L^0_K$,


поэтому, рассматривая линейные комбинации элементов класса,


мы не выйдем за его пределы. Это обстоятельство далее будет


использоваться без оговорок.





Как обычно, через $\Bbb N$ 


обозначается множество всех натуральных чисел $\{ 1,2,\dotsc \}$.


Решетка подполугрупп полугруппы натуральных чисел по сложению


обозначается через $\Sub (\Bbb N;+)$.





С точки зрения изучения свойств интервала $[L^0_K,F_K]$


бесконечные поля простой характеристики 


значительно отличаются от конечных и полей характеристики $0$.


Этому сложному случаю будет посвящена отдельная работа.





\section{Об одночленах замкнутых классов}





В данном параграфе будет доказана лемма, позволяющая, в частности, 


сделать вывод о том,


что для полей каждый из исследуемых классов порождается


своими одночленами. 





Везде далее через $K$ обозначаем поле.


Считаем, что в полиноме как сумме одночленов подобные 


члены приведены. 


Если $K=GF(p^n)$, то полагаем, что переменные входят в многочлены


в степенях меньших, чем $p^n$. 





Сформулируем без доказательства следующее простое


\medskip





{\bf Утверждение.} Пусть $d_1$ и $d_2$ --- различные натуральные числа


(меньшие, чем $p^n$, если $K=GF(p^n)$). Тогда существует ненулевое $a \in K$ 


такое, что $a^{d_1} \ne a^{d_2}$.


 


\begin{lem}


Каждый одночлен многочлена, лежащего в замкнутом классе, сам принадлежит


этому классу.


\end{lem}





{\bf Доказательство.} Пусть $C$ --- замкнутый класс, 


$f\in C$ и $f$ есть сумма


одночленов ${m_1},\dotsc,{m_l}$. Требуется доказать, что


$m_1,\dotsc,m_l\in C$.


Докажем лемму индукцией по $l$. При $l=1$ утверждение очевидно.


Пусть теперь оно выполнено для меньших, чем $l$, значений.


Обозначим через $d_i$ показатель переменной $x_1$ в одночлене $m_i$ 


(для всех $i\in \{1,\dotsc,l\}$).


В этой ситуации есть переменная, входящая в два одночлена с разными


степенными показателями, т.\,к. подобные одночлены приведены.


Не ограничивая общности, 


можно считать, что эта переменная есть $x_1$, что эти


одночлены суть $m_1$ и $m_2$, 


а соответствующие показатели удовлетворяют неравенству ${d_1}<{d_2}$.





В случае ${d_1}=0$ выполняется равенство


$m_1(0,\dotsc,x_r)=m_1(x_1,\dotsc,x_r)$.


Подставляя вместо ${x_1}$


константу $0\in {L^0_K}$, получим следующий многочлен из $C$:


$$


h=f(0,x_2,\dotsc,x_r)=m_1(x_1,\dotsc,x_r)


+\sum_{i=3}^l m_i(0,x_2,\dotsc,x_r).


$$


Так как ${d_2}>0$, то в $h\in C$ одночленов меньше, чем $l$.


Значит, по предположению индукции ${m_1}\in C$, $f-{m_1}\in C$,


и в многочлене $f-{m_1}=\sum\limits^l_{i=2} m_i$


число одночленов меньше, чем $l$. Теперь по предположению индукции


${m_2},\dotsc,{m_l}\in C$ и, наконец, ${m_1},{m_2},\dotsc,{m_l}\in C$.





Пусть теперь ${d_1}\ne 0$. 


Тогда по предыдущему утверждению 


существует такое ненулевое $a\in K$, 


что ${a^{d_1}}\ne {a^{d_2}}$.


Запишем одночлены в виде ${m_i}={x_1^{d_i}}{m_i^{\#}}({x_2},\dotsc,{x_r})$


и построим полином из $C$


$$


g=a^{d_1}f(a^{-1}x_1,x_2,\dotsc,x_r)-f(x_1,\dotsc,x_r)=


\sum_{i=1}^l (a^{d_1-d_i}-1){x_1^{d_i}}{m_i^{\#}}


=\sum_{i=2}^l (a^{d_1-d_i}-1){x_1^{d_i}}{m_i^{\#}}.


$$


Здесь $(a^{d_1-d_2}-1){x_1^{d_2}}{m_2^{\#}}\not\equiv 0$,


т.\,к. ${a^{d_1}}\ne {a^{d_2}}$.


Очевидно, $g$ --- сумма не более, чем $l-1$ одночленов. 


По предположению индукции $(a^{d_1-d_2}-1){m_2}\in C$, и поэтому


$m_2={(a^{d_1-d_2}-1)^{-1}}(a^{d_1-d_2}-1){m_2}\in C$.


Следовательно, $f-m_2\in C$,


и в $f-m_2$ не более $l-1$ одночленов. Аналогично случаю ${d_1}=0$ получаем


${m_1},{m_2},\dotsc,{m_l}\in C$.





Доказанная лемма дает основание утверждать, что каждый


замкнутый класс порождается своими одночленами.


Поэтому целесообразно выделить наиболее простые одночлены,


определяющие класс. 


Для полей характеристики $0$ это будет проделано в следующем параграфе. 





\section{Случай полей характеристики $0$}





В данном параграфе под $K$ понимается поле характеристики $0$,


и можно считать,


что $K$ содержит поле рациональных чисел $\Bbb Q$.


Как будет показано,


для таких полей интервал $[L^0_K,F_K^0]$ изоморфен 


решетке $\Sub (\Bbb N;+)$, 


а интервал $[L_K,F_K]$ прост (см. диаграмму 3).





\begin{picture}(400,100)


%% added 100 units to all X


\put(230,50){\oval(15,40)}





\put(230,30){\circle*{5}}


\put(230,70){\circle*{5}}


\put(250,40){\circle*{5}}


\put(250,80){\circle*{5}}





\put(230,20){$L^0_K$}


\put(252,78){$F_K$}


\put(221,77){$F^0_K$}


\put(250,30){$L_K$}





\put(230,30){\line(2,1){20}}


\put(230,70){\line(2,1){20}}


\put(250,40){\line(0,1){40}}





\put(130,45){$\Sub (\Bbb N;+)$}


\put(183,48){\vector(1,0){47}}


\put(200,50){$\cong$}





\put(180,5){Диаграмма 3}


\end{picture}





Как показано в предыдущем параграфе, 


каждый замкнутый класс порождается своими одночленами.


Oпределим достаточно простые множества


таких одночленов.


В связи с этим введем несколько понятий.





Одночлен вида $x_1\dotsc x_n$, являющийся произведением $n$ 


различных переменных,


будем называть простым ($1$ также считаем простым одночленом).





\begin{lem}


Если замкнутый класс содержит одночлен степени $D$,


то он содержит и простой одночлен той же степени.


\end{lem}





{\bf Доказательство.}


Возьмем произвольный одночлен $m$ из замкнутого класса $C$.


Не ограничивая общности, можно считать,


что $m={x_1^d m_1(x_2,\dotsc,x_r)}$, $d>0$. 


Подставив вместо $x_1$ сумму из $d$ 


переменных $y_1+\dotsb +y_d\in L^0_K$, получим


$$
(y_1+\dotsb +y_d)^d m_1=({k \cdot y_1\dotsc y_d} +g)m_1\in C,


$$


где $k\in \Bbb N$ и многочлен $g$ не содержит одночленов,


подобных $y_1\dotsc y_d$.


По лемме~1 имеем ${k y_1\dotsc y_d}m_1\in C$.


Подставляя вместо переменной $y_1$ 


одночлен $k^{-1}y_1\in L^0_K\subset C$ 


($k^{-1}\in \Bbb Q\subset K$),


получим 


$$


{y_1\dotsc y_d}m_1(x_2,\dotsc,x_r)\in C.


$$


Применяя к переменным $x_{2},\dotsc,x_r$ в полученном одночлене рассуждения,


аналогичные проведенным для $x_1$, получим


$$


{y_1\dotsc y_D}\in C.\quad\square


$$





Отождествляя подходящим образом переменные в простом одночлене,


с помощью итеративных операций легко получить любые одночлены той


же степени.


Поэтому доказанная лемма дает основание утверждать, 


что любой замкнутый класс порождается своими {\em простыми} одночленами.





\begin{cor}


Интервал $[L_K,F_K]$ прост.


\end{cor}





{\bf Доказательство.}


Пусть $C$ --- замкнутый класс из $[L_K,F_K]$, отличный от $L_K$.


Покажем, что $C=F_K$.


Класс $C$ содержит одночлен степени $l>1$, т.\,к. $C\ne L_K$.


По лемме 2 он содержит одночлен $x_1\dotsc x_l$.


Подставляя вместо $l-2$ последних переменных 


$1\in L_K$,


получим ${x_1 x_2}\in C$.


Заметим, что 


\begin{align*}


&{x_1x_2}*{x_1x_2}={x_1x_2x_3},\\


&{x_1x_2x_3}*{x_1x_2}={x_1x_2x_3x_4},\\


&\dotsc \\


&{x_1x_2\dotsc x_{n-1}}*{x_1x_2}={x_1x_2x_3\dotsc x_n},\\


&\dotsc 


\end{align*}


где $*$ --- итеративная операция подстановки (см. {\S\,1}). 


Класс $C$ содержит все простые одночлены.


Значит, $C=F_K$.





Итак, интервал $[L_K,F_K]$ прост. Oсталось выяснить


устройство интервала $[L^0_K,F^0_K]$.


До конца данного параграфа слова ``замкнутый класс'' 


будут обозначать замкнутый класс из $[L^0_K,F^0_K]$.





Назовем длиной одночлена его степень без единицы.





\begin{lem}


Положительные длины одночленов из замкнутого класса


суть в точности элементы подполугруппы из $\Bbb N$,


порожденной длинами одночленов 


из произвольного порождающего множества этого замкнутого класса.


\end{lem}





{\bf Доказательство.}


Рассмотрим замкнутый класс $C=\Class{ {f_1},\dotsc,{f_k},\dotsc}$,


где $f_1,\dotsc,f_k,\dotsc$ --- одночлены с ненулевыми длинами


$l_1,\dotsc,l_k,\dotsc$, и пусть $S=\Close{l_1,\dotsc,l_k,\dotsc}$ ---


подполугруппа, порожденная этими длинами.


Класс $C$ является подалгеброй, 


порожденной множеством $L^0_K\cup\{f_1,\dotsc,f_k,\dotsc \}$


с помощью операций $\zeta$, $\tau$, $\Delta$, $\nabla$, $*$.


Докажем индукцией по построению этой подалгебры,


что длины всех одночленов из $С$ лежат в $S$.


База индукции очевидна.


Переходя к шагу индукции, отметим,


что операции $\zeta$, $\tau$, $\Delta$, $\nabla$


не изменяют степени одночленов.


Поэтому достаточно рассмотреть действие на степени


операции $*$.





При подстановке $m*f$, 


где $m=x_1^{\alpha}m^{\#}(x_2,\dotsc,x_s)$ ---


одночлен степени $1+q$, $f=m_{1}+\dotsb +m_{r}$, $m_{i}$ --- одночлены,


$\deg m_{i}=1+{q_i}$, $q,q_1,\dotsc,q_r\in S$, 


получаются одночлены с длинами вида 


${k_0 q}+{k_1 q_1}+\dotsc +{k_l q_r}\in S$.


В самом деле,


$$


(m_{1}+\dotsc +m_{r})^{\alpha}=


\sum_{t=0}^T a_t m_{1}^{\beta _{1t}}\dotsc m_{r}^{\beta {r t}},


$$


где $\beta _{1t}+\dotsb +\beta _{r t}=\alpha$


для любого $t\in\{0,\dotsc,T\}$.


Если $a_t\ne 0$, то слагаемое с номером $t$ имеет степень


$$


(1+q_1)\beta _{1t}+\dotsb +(1+q_r)\beta _{r t}=


\sum_{j=1}^{r}\beta _{jt}+\sum_{j=1}^{r}\beta _{jt}q_j=


\alpha+q_1\beta _{1t}+\dotsb +q_r\beta _{r t}.


$$


При подстановке $m*(m_{1}+\dotsc +m_{r})$ получим


$$


(m_{1}+\dotsb +m_{r})^{\alpha}m^{\#}=


\sum_{t=0}^T a_t m_{1}^{\beta _{1t}}\dotsc


m_{r}^{\beta {r t}}m^{\#}.


$$


Здесь слагаемое с номером $t$ имеет степень


$$


\alpha +q_1\beta _{1t}+\dotsb +q_r\beta _{r t}+(1+q-\alpha)=


1+q_1\beta _{1t}+\dotsb +q_r\beta _{r t}+q,


$$


т.\,к. $\deg m^{\#} =1+q-\alpha$.


Таким образом, длины всех вновь полученных одночленов


лежат в $S$.





Общий случай сводится


к предыдущему, т.\,к.


$\Delta(M_1+\dotsb +M_s) = \Delta M_1+\dotsb +\Delta M_s$,


$(M_1+\dotsb +M_s)*f = M_1*f+\dotsb +M_s*f$.





Чтобы доказать, что положительные длины одночленов из $C$


принимают {\em все} значения из $S$,


рассмотрим простые одночлены с длинами $l_1,\dotsc,l_k,\dotsc$


По лемме 2 такие одночлены в $C$ есть.


Заметим теперь, что длины простых одночленов при подстановке


складываются


$$


x_1\dotsc x_{L_1 +1} * x_1\dotsc x_{L_2 +1}=x_1\dotsc x_{L_1+L_2+1}.


$$


Это, в частности, означает, что длины простых одночленов


из $C$ пробегают всю полугруппу $S$.





Основной результат данной работы дает


\medskip





{\bf Теорема. }{\it


Интервал $[L^0_K,F^0_K]$ изоморфен решетке $\Sub (\Bbb N;+)$.


}


\medskip





{\bf Доказательство.}


Каждый замкнутый класс порожден своими простыми одночленами.


Поставим в соответствие каждой подполугруппе $S\in \Sub (\Bbb N;+)$


замкнутый класс $\phi (S)$, порожденный всеми простыми одночленами,


длины которых принимают значения из $S$.


По лемме~3 положительные длины одночленов из замкнутого класса


суть в точности элементы подполугруппы из $\Bbb N$


порожденной


длинами одночленов из порождающего множества этого замкнутого класса.


В силу леммы~2 эти длины задают длины простых одночленов


замкнутого класса, которые в свою очередь определяют сам класс.


Таким образом, построенное соответствие


есть взаимно однозначное отображение на $[L^0_K,F^0_K]$.


Легко видеть, что оно сохраняет все включения: 


$$


\forall S_1,S_2\in \Sub (\Bbb N;+)\,\,


( S_1\subset S_2 \Leftrightarrow 


\phi (S_1)\subset \phi (S_2)).\quad\square


$$





\begin{cor}


Любой класс из $[L^0_K,F_K]$ конечнопорожден.


\end{cor}





{\bf Доказательство.}
Очевидно, $L_K=\Class{1}$ и $F_K=\Class{1, {x_1 x_2} }$ конечнопорождены. 


Известно также, 


что любая полугруппа из $\Sub (\Bbb N;+)$ конечнопорождена.


В качестве порождающего множества для класса из $[L^0_K,F^0_K]$


возьмем множество всех тех простых одночленов этого класса, 


длины которых принимают все значения 


из произвольного {\em конечного} порождающего множества


соответствующей ему полугруппы. Длины таких одночленов при подстановке


складываются и порождают таким образом исходную подполугруппу, которая


и определяет класс. Значит, классы из $[L^0_K,F^0_K]$


также конечнопорождены.





\section{Заключительные замечания}





Первые результаты о замкнутых классах полиномов


были получены еще в 50-е годы, в частности, С.В.\,Яблонским [4]


и А.В.\,Кузнецовым [5].


С тех пор этой теме был посвящен ряд работ, продолжением которого


является и данная работа.


Несколько необычной для данной тематики является бесконечность


полей, над которыми рассматриваются полиномы.





В случае полей характеристики $0$ удалось


описать интервал $[L^0_K,F_K]$ с точностью до строения


решетки $\Sub (\Bbb N;+)$, 


хорошо известной в теории полугрупп.


Ее строением еще в 60-е годы заинтересовались


в ходе изучения свойств полугрупп


с решетками подполугрупп, удовлетворяющими нетривиальному тождеству.


Вопрос о том, 


удовлетворяет ли $\Sub (\Bbb N;+)$ какому-либо нетривиальному тождеству,


был поставлен Л.Н.\,Шевриным 


([6], задача 2.74).


Задача была решена в [7].


В статье [8] было отмечено, 


что конечные подрешетки из $\Sub (\Bbb N;+)$ ---


это в точности все конечные ограниченные снизу %%%(lower bounded)


решетки.


Полное доказательство этого факта приведено в ([9],


теорема~28.1).


Также важным для нас свойством решетки $\Sub (\Bbb N;+)$


является ее счетность.








В заключение автор выражает благодарность


Е.В.\,Суханову, Л.Н.\,Шеврину и


всем, кто высказал свои замечания по данной работе.








\begin{thebibliography}{9}





\bibitem{1}


Семигродских А.П., Суханов Е.В. {\em О замкнутых классах полиномов


над конечным полем} // Дискрет. матем. -- 1997. -- Т.\,9. -- \No\,4. -- 


C.\,50--62.


\bibitem{2}


Мальцев А.И. {\em Итеративные алгебры Поста}. --


Новосибирск: Изд-во Новосибирск. ун-та, 1974. -- 80 с.


\bibitem{3}


Крохин А.А., Сафин К.Л., Суханов Е.В.


{\em О строении решетки замкнутых классов полиномов} // Дискрет. матем. --


1997. -- Т.\,9. -- \No\,2. -- C.\,24--39.


\bibitem{4}


Яблонский С.В. {\em Функциональные построения в $k$-значной логике} //


Тр. Матем. ин-та АН СССР. -- М., 1958. -- Т.\,51. -- C.\,5--142.


\bibitem{5}


Кузнецов А.В. {\em О бесповторных контактных схемах и бесповторных


суперпозициях функций алгебры логики} //


Тр. Матем. ин-та АН СССР. -- М., 1958. -- Т.\,51. -- C.\,186--225.


\bibitem{6}


{\em Свердловская тетрадь. Нерешенные задачи теории полугрупп}. --


Свердловск:


Уральск. ун-т, 1979. -- Вып.\,3. -- 41 с.


\bibitem{7}


Репницкий В.Б., Кацман С.И.


{\em Коммутативные полугруппы, решетка подполугрупп которых


удовлетворяет нетривиальному тождеству} // Матем. сб. -- 


1988. -- Т.\,137. -- \No\,4. -- C.\,462--482.


\bibitem{8}


Repnitski\u\i{} V.B. 


{\em On finite lattices which are embeddable in subsemigroup lattices}


// Semigroup Forum. -- 1993. -- V.\,46. -- P.\,388--397.


\bibitem{9}


Shevrin L.N., Ovsyannikov A.J.


{\em Semigroups and their subsemigroup lattices}. --


Dordrecht--Boston--London: Kluwer Acad. Publ., 1996. -- 390 p.





\end{thebibliography}





\vskip 2cm





\post{\it Уральский государственный университет}{\it Поступила}


\post{}{09.01.1998}





\label{end}


\end{document}





